
Le tableau conserve en outre les réalisations quantiques que sa classification
abrège. Les trois leptons chargés et les six saveurs de quarks appartiennent à la
complétion extérieure : CAR conduit à Pauli et à Fermi--Dirac, sans attribuer ces
relations à un TRIT isolé. Les quarks occupent trente-six cellules lorsque l'on
conserve couleur et antiparticule avant le quotient. Le support neutre commun
\(\mathcal H_\nu=\bigoplus_{j=1}^{4}\mathcal H_{\nu,G_j}\) a pour dimensions
\(2+4+6+8=20\). Les symboles \(G_j\) désignent des strates de profondeur de l'atlas ;
\(f_e,f_\mu,f_\tau\) désignent la base d'interaction. En particulier, \(G_4\) est
une strate du support neutre commun et ne s'identifie pas à un état stérile ou
exotique sans projecteur supplémentaire. PMNS transporte la base d'interaction vers la
base des vecteurs propres ; il ne crée ni valeurs propres, ni échelle de masse, ni sélecteur de
profondeur. Le photon conserve aller et retour dans la
carte de jauge nulle : \(m=0\) ne signifie pas évaluer \(\chi_0\) avec un exposant très
négatif. \(W^\pm\) et \(Z^0\) sont des représentations bosoniques massives, avec
conjugaison de charge pour \(W^+\leftrightarrow W^-\). Enfin,
l'antisection s'obtient par inversion et renversement de bande ; la masse est paire
sous cette involution et n'est pas doublée comme second paramètre indépendant.

Tous les secteurs conservent le préfixe causal
\[
\text{extension survivante}\longrightarrow\text{route observable}
\longrightarrow\text{registre chronologique enrichi}.
\]
Le lecteur change ensuite. Dans le secteur matériel positif, la chaîne se poursuit par
\(\text{signature}\to\text{caractère}\to\text{tours}\to\text{carte dimensionnelle}\).
Le secteur nul évite le caractère positif et fixe \(m=0\) ; les hadrons composent
d'abord des enregistrements chronologiques enrichis ; les secteurs topologiques utilisent des invariants de
fusion ou de tresse ; et la virtualité se reconnaît à l'obstruction terminale. Le
domaine de ces applications est la classification complète précédente. Un essai
numérique sur un sous-ensemble n'établit que le cas fini testé et ne
réduit pas ce domaine.

Les dénombrements qui accompagnent ce tableau relèvent de quotients différents :
\[
81=25_{\ell^-}+20_\nu+20_d+16_u,
\]
\[
56\ \text{familles à fibres}\quad
41\times1+10\times2+2\times3+1\times4+2\times5=81,
\]
\[
13\ \text{multisections},\qquad324=81\times4\ \text{généalogies}.
\]
De même, les \(192\)-A classes de représentation associées à
\(\operatorname{diag}(2,6,2,2,4)\) et l'indice réticulaire \(192\)-B de
\(\operatorname{SNF}(2,2,2,2,12)\), avec \(52\) signatures brutes, ne sont pas le
Le cardinal du catalogue historique d'exposition ne compte pas les particules individuelles et ne définit pas la taxonomie matérielle.
non des particules individuelles.
